\section{The OP5 comparison model}\label{sec:op5}

OP5 is the question of the cosmological constant and the vacuum balance.  In its exact form it is open; only a local detector
version is closed in the Companion.  Since late October 2026 the work on it has turned into a precise mathematical question that
can be answered independently of the physical reading: does a fixed system of field equations, built from the probe--closure
system of Section~\ref{sec:octonions}, have solutions on an open region?  This section reports the answer so far.  Everything here
concerns one declared \emph{comparison model}, not the framework as a whole, and none of it closes OP5.

\subsection{Three exclusions}

Before the model was fixed, three general statements were established (layer v4.9.0-r69).  A change of frame cannot remove a defect
in the vacuum balance.  None of the routes considered fixes the absolute vacuum term: state, frame and normalization leave it free, so
\(\Lambda\) can at best appear as an integration constant.  And under torsion-freeness and the ordinary Einstein vacuum equation the
geometric cosmological constant must vanish wherever the cross current of the probes is non-zero --- an exclusion under stated
assumptions.

\subsection{The model}

The unknown is a spinor field \(\Psi\colon\mathbb R^4\to\mathbb C^{12}\).  With
\(D_\mu=\partial_\mu+iA_\mu\), the coframe \(E^a{}_\mu=\operatorname{Re}(\Psi^\dagger Q_aD_\mu\Psi)\) and its inverse \(V=E^{-1}\), the equations are
\[
F:=i\,V^\mu{}_a K_a D_\mu\Psi-h\,(Y\otimes\beta)\Psi=0,\qquad
g:=4-\tfrac h2\,\Psi^\dagger(Z\otimes\beta)\Psi=0,
\]
completed to \(R=(F,G,dg)\) with \(24+4+4\) real components, where \(K_a=I_3\otimes(I,\alpha_i)\), \(Q_a=S\otimes(I,\alpha_i)\), the connection is
\(A_0=0\), \(A_k=G(x)\,\mathrm{spin}_k(x)\otimes I_4\), and the background is \(h=h(t,r)\), \(A=M(t)\,d\varphi\).  The geometry is thus
reconstructed from the field itself.  The model is a declared classical comparison sector: its couplings and background are chosen,
not derived from the kernel.  A germ at \((t,x,y,z)=(1,\tfrac35,0,0)\) is known to \(70\) digits up to fifth order and satisfies \(R\) to
third order.

The method is the formal theory of partial differential equations~\cite{bryant,seiler}.  One computes the \emph{symbol} \(g_p\) of the
system at order \(p\), tests whether it is \emph{involutive} (Cartan's test), and checks \emph{formal integrability}, i.e.\ whether
prolongation produces new conditions of lower order.  An involutive, formally integrable analytic system has local analytic solutions,
and its Cartan characters say how many free functions of how many variables they depend on.

\subsection{An exact gauge symmetry}

The degree-4 symbol has rank \(528\) and a left kernel of dimension \(112\), and the additional third-order condition \(\Omega\) on the
germ is unique.  Its characteristic set consists of a light cone \(q\) and two skew planes.  One direction of the solution space is
free, and it is not an accident (layer v4.9.0-r71):

\begin{proposition}[Gauge symmetry]\label{prop:gauge}
For the background \(h=h(t,r)\), \(A=M(t)\,d\varphi\), the variation \(\delta\Psi=f\,(\rho\partial_z-z\partial_\rho)\cdot D\Psi\) with an
arbitrary function \(f\) is an exact symmetry of the system, and \(\Omega\) is invariant under it.
\end{proposition}

Two pre-registered hypotheses about this direction were refuted: a hidden condition at second order off the equator (H49), and a
back-reaction of \(\Omega\) on \(\Psi\) through it (H50).  The symmetry has to be fixed before an existence theorem can apply.

\subsection{Cartan test of the gauged system}

Fixing the gauge algebraically, every transversal choice gives the same symbol (layer v4.9.0-r72):
\[
\dim g'_p=117,\ 188,\ 273,\ 371,\ 481,\ 606,\ 745\qquad(p=2,\dots,8).
\]
The pre-registered test H51 predicted from the growth of these numbers that \(g'\) becomes involutive first at \(p=6\); this was
confirmed in three coordinate systems, with Cartan characters \(481,111,14,0\).  No free function of four variables is left; formally the
solutions depend on \(14\) functions of three variables, i.e.\ on data on a hypersurface.

\subsection{Four hidden conditions}

The gauged system is involutive but not formally integrable at order six.  Prolonging to degree eight reduces the dimension of the
sixth-order fibre from \(852\) to \(848\): there are four additional conditions on 6-jets (layers v4.9.0-r72, r73).  Their status:
\begin{itemize}
\item The Spencer cohomology of \(g'\) allows at most four such conditions (\(H^{4,2}=0\), \(H^{5,2}=4\)), and a Weyl bound on the measured
error norms excludes that the drop is numerical noise.
\item They were found again at a second germ, \(50\%\) away, with error bounds \(710\)- and \(1390\)-fold (pre-registered test H55,
confirmed).
\item They are genuine sixth-order equations.  They act only on light-like plane waves, one polarization per null direction, in the
form \(d_k(\xi)\,s(\xi,u)\) with the four quadrics \(d_k\) that vanish on the two special planes: they are the curvature remainders of the
four relations \(q\,d_k\,\Omega\) found at the level of the symbol.
\end{itemize}

\subsection{The extended system and a precision limit}

Adding the four conditions gives an extended symbol \(g''\) with
\(\dim g''_p=477,594,723,865\) (\(p=6,\dots,9\)); by all evidence it is involutive first at \(p=8\), with characters \(723,130,12,0\), i.e.\
\(12\) functions of three variables (layer v4.9.0-r74).  The corresponding pre-registered test was formally invalid, because five rank
decisions fell just short of the required singular-value gap.  The Spencer cohomology of \(g''\) leaves room for up to four further
hidden conditions on 7-jets and three on 8-jets.

Deciding them meets a structural precision limit.  The symbol of the field equation alone is well conditioned, but \(\Omega\) and the
four conditions act only on the cone and the planes and are, as rows, nearly dependent on it (independent weight about \(10^{-8}\)).
Double precision separates one or two such levels, not three.  A staged elimination saves memory but does not lift the limit, and a
reformulation by plane waves on the cone does not span the symbol spaces.  A recomputation of the whole chain at \(100\) digits has
started.

\subsection{Larger orders on parallel runners}

As a test of distributed computation, the symbol dimensions of the frozen \(\Omega\)-symbol were computed on parallel runners, one order
per job, up to \(p=12\) (matrices up to \(12\,323\times10\,920\)).  All known values were reproduced, and the new ones,
\(\dim g_{11},\dim g_{12}=1610,1896\) and \(\dim g'_{10},\dim g'_{11},\dim g'_{12}=1065,1246,1441\), have constant second differences
\(14\) for the gauged symbol from \(p=7\) on, as the \(14\) functions of three variables require.  This is an exploration on the older
germ, not a registered test.

\medskip\noindent\emph{Status.}  For the comparison model: an exact gauge symmetry; an involutive gauged symbol at order six (confirmed);
four hidden sixth-order conditions (confirmed at two points); an extended symbol involutive at order eight by all evidence but not
formally confirmed; further conditions not excluded.  If the extended system turns out involutive and formally integrable, the
Cartan--K\"ahler theorem would give local analytic solutions depending on \(12\) functions of three variables --- for this model only.
That would be a local existence statement, not a closure of OP5 and not a property of the framework.
